%----------
%	SKILLS, PROMPTS Y HERRAMIENTAS
%----------

Toda la metodología del sistema vive en ficheros \texttt{SKILL.md}: un documento
por unidad de conocimiento, con un encabezado (nombre y descripción) y el cuerpo
con las instrucciones. La descripción es lo que el agente ve antes de decidir si
carga la skill, y por tanto lo que gobierna el enrutado. El cuerpo solo se lee
cuando hace falta.

\subsection{Tipos de skill}

Hay tres tipos, con papeles distintos:

\begin{itemize}
    \item \textbf{Skills de variable} (cinco). Contienen la metodología paso a
    paso de cada variable. Se generan automáticamente desde
    \texttt{variables.json}, de modo que el codebook y las skills nunca se
    desincronizan. Cada agente carga la suya.
    \item \textbf{Skills auxiliares} (tres): \texttt{guia\_regla\_inversion},
    \texttt{guia\_lenguaje\_inclusivo} y \texttt{verificar\_evidencias}.
    Cualquier agente las carga bajo demanda cuando duda.
    \item \textbf{Resúmenes de guías}. Condensaciones de las guías reales de
    lenguaje no sexista (Sección~\ref{rag}), disponibles como orientación rápida.
\end{itemize}

Que las skills de variable se generen desde \texttt{variables.json} tiene una
consecuencia importante para la validez del experimento: el nivel B0 y el nivel
B1 parten del mismo fichero fuente. El \textit{baseline} recibe ese contenido
inyectado en el \textit{prompt}; el agente lo carga como skill. El contenido es
byte a byte el mismo, así que la comparación aísla el mecanismo de entrega y no
un cambio en la metodología.

\subsection{Herramientas}

Las herramientas del agente son deliberadamente pocas y de propósito claro:

\begin{itemize}
    \item \texttt{list\_skills} devuelve solo los metadatos (nombre y
    descripción) de las skills visibles. Es lo que hace efectiva la
    \textit{progressive disclosure}: el agente elige a ciegas del cuerpo.
    \item \texttt{read\_skill} entrega el cuerpo completo de una skill concreta.
    \item \texttt{verificar\_evidencias} comprueba, contra el mismo texto que vio
    el modelo, qué fragmentos citados son literales.
\end{itemize}

\subsection{Salida estructurada y verificación de evidencias}

Cada agente cierra con un objeto JSON de tres campos (\texttt{codigo},
\texttt{explicacion}, \texttt{evidencias}), descrito en la
Sección~\ref{sec:zeroshot_iris}. Sobre las evidencias se aplica una comprobación
automática: solo se conservan los fragmentos que aparecen \emph{literalmente} en
el texto original. Esta verificación cumple dos funciones. Refuerza la
trazabilidad exigida por el marco \textit{human-in-the-loop}, y evita un fallo
sutil que arrastraban implementaciones anteriores, donde el texto se alteraba
antes de contrastar las evidencias y estas se descartaban por error. Aquí la
comprobación se hace contra el texto exacto que recibió el modelo.
